\documentclass[aspectratio=43]{beamer}

\usetheme{Madrid}
\usecolortheme{default}
\usefonttheme{professionalfonts}
\setbeamertemplate{navigation symbols}{}
\setbeamertemplate{blocks}[rounded][shadow=true]

\usepackage{amsmath,amssymb,graphicx,booktabs,array}
\usepackage{hyperref}
\usepackage{listings}
\usepackage{xcolor}
\usepackage{tikz}

\graphicspath{{plots/}{./}}

% ---- Custom color palette ----
\definecolor{Midnight}{HTML}{0F172A}
\definecolor{Slate}{HTML}{1E293B}
\definecolor{ElectricBlue}{HTML}{38BDF8}
\definecolor{Violet}{HTML}{8B5CF6}
\definecolor{Rose}{HTML}{F43F5E}
\definecolor{SoftWhite}{HTML}{F8FAFC}
\definecolor{CoolGray}{HTML}{CBD5E1}

% ---- Global beamer styling ----
\setbeamercolor{background canvas}{bg=SoftWhite}
\setbeamercolor{normal text}{fg=Midnight,bg=SoftWhite}
\setbeamercolor{frametitle}{fg=SoftWhite,bg=Midnight}
\setbeamercolor{title}{fg=SoftWhite}
\setbeamercolor{subtitle}{fg=CoolGray}
\setbeamercolor{author}{fg=SoftWhite}
\setbeamercolor{institute}{fg=CoolGray}
\setbeamercolor{date}{fg=CoolGray}
\setbeamercolor{palette primary}{fg=SoftWhite,bg=Midnight}
\setbeamercolor{palette secondary}{fg=SoftWhite,bg=Slate}
\setbeamercolor{palette tertiary}{fg=SoftWhite,bg=Violet}
\setbeamercolor{palette quaternary}{fg=SoftWhite,bg=ElectricBlue}
\setbeamercolor{section in toc}{fg=Violet}
\setbeamercolor{subsection in toc}{fg=Slate}
\setbeamercolor{itemize item}{fg=Violet}
\setbeamercolor{itemize subitem}{fg=ElectricBlue}
\setbeamercolor{enumerate item}{fg=Violet}
\setbeamercolor{block title}{fg=SoftWhite,bg=Violet}
\setbeamercolor{block body}{fg=Midnight,bg=violet!8}
\setbeamercolor{alerted text}{fg=Rose}
\setbeamercolor{caption name}{fg=Violet}

\setbeamerfont{title}{series=\bfseries,size=\LARGE}
\setbeamerfont{frametitle}{series=\bfseries,size=\large}
\setbeamerfont{block title}{series=\bfseries}

\hypersetup{
  colorlinks=true,
  linkcolor=Violet,
  urlcolor=ElectricBlue,
  citecolor=Rose
}

\lstset{
  basicstyle=\ttfamily\scriptsize,
  breaklines=true,
  frame=single,
  columns=fullflexible,
  showstringspaces=false,
  keywordstyle=\color{Violet}\bfseries,
  commentstyle=\color{Slate},
  stringstyle=\color{Rose},
  rulecolor=\color{ElectricBlue},
  backgroundcolor=\color{SoftWhite}
}

\setbeamertemplate{title page}{
  \begin{tikzpicture}[remember picture,overlay]
    \fill[Midnight] (current page.south west) rectangle (current page.north east);
    \fill[Violet,opacity=0.18] ([xshift=-1.0cm,yshift=1.5cm]current page.south east) circle (3.0cm);
    \fill[ElectricBlue,opacity=0.14] ([xshift=1.0cm,yshift=-1.0cm]current page.north west) circle (2.8cm);
  \end{tikzpicture}
  \vbox{}
  \vfill
  \begin{center}
    {\usebeamerfont{title}\usebeamercolor[fg]{title}\inserttitle\par}
    \vskip0.8em
    {\usebeamerfont{subtitle}\usebeamercolor[fg]{subtitle}\insertsubtitle\par}
    \vskip1.5em
    {\usebeamerfont{author}\usebeamercolor[fg]{author}\insertauthor\par}
    \vskip0.4em
    {\usebeamerfont{institute}\usebeamercolor[fg]{institute}\insertinstitute\par}
    \vskip1.0em
    {\usebeamerfont{date}\usebeamercolor[fg]{date}\insertdate\par}
  \end{center}
  \vfill
}

\title[P vs NP]{A Novel Perspective on the P vs NP Problem\\with an AI Co-Pilot Exploration Framework}
\author[Samuel Castillo]{Samuel Castillo}
\institute[UMass Amherst]{University of Massachusetts Amherst}
\date{April 4, 2026}


\begin{document}

\begin{frame}[plain]
\titlepage
\end{frame}

\begin{frame}{Contents}
\tableofcontents
\end{frame}

\section{Motivation}

\begin{frame}{Why This Problem Matters}
\begin{block}{Core idea}
\begin{itemize}
    \item The P vs\@. NP problem asks whether every efficiently verifiable problem is also efficiently solvable.
    \item It sits at the center of complexity theory, optimization, cryptography, and algorithm design.
    \item The practical issue is \textbf{computational feasibility}: which problems stay solvable under real limits on time, memory, energy, and hardware.
    \item This presentation reframes the paper around \textbf{Time-to-Solve (TTS)}, large-scale systems, and an \textbf{AI-assisted exploration framework}.
\end{itemize}
\end{block}
\end{frame}

\begin{frame}{Real-World Motivation: Travel Systems and Optimization}
\begin{itemize}
    \item Large travel systems must forecast demand, allocate inventory, schedule under disruptions, and optimize itineraries.
    \item Many of these tasks resemble classical combinatorial optimization problems.
    \item The canonical example is the \textbf{Traveling Salesman Problem (TSP)}:
\end{itemize}

\medskip
\centering
\emph{Given a set of locations and pairwise distances, what is the shortest route that visits each location once and returns to the origin?}

\medskip
\raggedright
\begin{itemize}
    \item Exact solutions typically become infeasible at scale.
    \item Practical systems rely on heuristics, approximations, decomposition, and predictive models.
\end{itemize}
\end{frame}

\section{Formal Definitions}

\begin{frame}{Class $\mathsf{P}$}
\[
\mathsf{P}
=
\left\{
L \subseteq \Sigma^*
\;\middle|\;
\exists \text{ deterministic Turing machine } M,
\exists k \in \mathbb{N},
\text{ such that } M \text{ decides } L
\text{ in time } O(n^k)
\right\}
\]

\begin{itemize}
    \item Informally: problems solvable in polynomial time.
    \item Polynomial time is usually treated as the mathematical model of tractability.
\end{itemize}
\end{frame}

\begin{frame}{Class $\mathsf{NP}$}
\[
\mathsf{NP}
=
\left\{
L \subseteq \Sigma^*
\;\middle|\;
\exists \text{ nondeterministic Turing machine } M,
\exists k \in \mathbb{N},
\text{ such that } M \text{ decides } L
\text{ in time } O(n^k)
\right\}
\]

\medskip
Equivalent verifier-based definition:
\[
x \in L \iff \exists y \in \Sigma^*,\ |y| \le p(|x|),\ V(x,y)=1
\]
where $V$ is a deterministic polynomial-time verifier.
\end{frame}

\begin{frame}{The Central Question}
\[
\text{Does } \mathsf{P} = \mathsf{NP}\ ?
\]

\medskip
\begin{block}{Equivalent wording}
Does every language that admits a polynomial-time verifier also admit a polynomial-time decision algorithm?
\end{block}

\begin{itemize}
    \item If \(\mathsf{P}=\mathsf{NP}\), many currently intractable search problems would become polynomial-time solvable.
    \item If \(\mathsf{P}\neq\mathsf{NP}\), then the gap between verification and search is fundamental.
\end{itemize}
\end{frame}

\section{Feasibility and Runtime Growth}

\begin{frame}{Time-to-Solve (TTS) and Feasibility}
\begin{itemize}
    \item Complexity is not just a theoretical label; it is a practical limit on resources.
    \item Feasibility depends on:
    \begin{itemize}
        \item time,
        \item memory,
        \item energy and hardware,
        \item scalability as input size grows.
    \end{itemize}
    \item Polynomial growth is usually manageable.
    \item Exponential growth quickly outruns any realistic compute budget.
\end{itemize}
\end{frame}

\begin{frame}{Scaling Behavior by Complexity Class}
\begin{itemize}
    \item \textbf{$O(1)$}: constant time
    \item \textbf{$O(\log n)$}: logarithmic time
    \item \textbf{$O(n)$}: linear time
    \item \textbf{$O(n^2)$}: quadratic time
    \item \textbf{$O(2^n)$}: exponential time
\end{itemize}

\medskip
Key idea: the difference between polynomial and exponential time is not a small efficiency gap. It is often the difference between \textbf{practical} and \textbf{impossible}.
\end{frame}

\begin{frame}{Real Runtime Comparison}
Assume a machine performing approximately $10^9$ basic operations per second.

\small
\begin{center}
\begin{tabular}{@{}cccc@{}}
\toprule
Input $n$ & $O(n)$ & $O(n^2)$ & $O(2^n)$ \\
\midrule
10  & $10^{-8}$ s & $10^{-7}$ s & $1.0\times10^{-6}$ s \\
30  & $3\times10^{-8}$ s & $9\times10^{-7}$ s & $1.07$ s \\
50  & $5\times10^{-8}$ s & $2.5\times10^{-6}$ s & $\approx 13$ days \\
70  & $7\times10^{-8}$ s & $4.9\times10^{-6}$ s & $\approx 37{,}000$ years \\
90  & $9\times10^{-8}$ s & $8.1\times10^{-6}$ s & $\approx 39$ billion years \\
100 & $10^{-7}$ s & $10^{-5}$ s & $\approx 40$ quintillion years \\
\bottomrule
\end{tabular}
\end{center}

\normalsize
By modest input sizes, exponential-time algorithms already exceed any realistic time horizon.
\end{frame}

\begin{frame}{Interpretation in TTS Terms}
\begin{block}{Reframing the conjecture}
Are problems that appear to require exponential time actually solvable in polynomial time?
\end{block}

\begin{itemize}
    \item Verification can stay easy while search explodes combinatorially.
    \item That asymmetry is the practical face of the P vs\@. NP problem.
    \item Better hardware helps, but exponential growth still wins.
\end{itemize}
\end{frame}

\section{Intuition and NP-Completeness}

\begin{frame}{Intuition: Verification vs. Search}
A useful mental model:
\begin{itemize}
    \item \textbf{Verifying} a candidate solution may be fast.
    \item \textbf{Finding} that solution may require massive search.
\end{itemize}

\medskip
Example: \textbf{Sudoku}
\begin{itemize}
    \item Checking a completed grid is quick.
    \item Constructing a valid solution can require extensive search.
\end{itemize}
\end{frame}

\begin{frame}{NP-Complete Problems}
A problem is NP-complete if:
\begin{itemize}
    \item it belongs to $\mathsf{NP}$, and
    \item every problem in $\mathsf{NP}$ reduces to it in polynomial time.
\end{itemize}

\medskip
Examples:
\begin{itemize}
    \item Traveling Salesman Problem
    \item Subset Sum
    \item Knapsack
\end{itemize}

\medskip
A polynomial-time algorithm for any NP-complete problem would imply
\[
\mathsf{P}=\mathsf{NP}.
\]
\end{frame}

\section{Computational Perspective}

\begin{frame}{Empirical Runtime Experiments}
Even though the conjecture is theoretical, experiments help build intuition.

\begin{lstlisting}[language=C]
for (int n = 10; n <= 1000; n *= 2) {
    auto start = now();
    solve_instance(n);
    auto end = now();
    log(n, end - start);
}
\end{lstlisting}

\begin{itemize}
    \item Runtime plots often make the polynomial vs. exponential gap obvious.
    \item Empirical scaling does not prove complexity class membership, but it helps expose growth behavior.
\end{itemize}
\end{frame}

\begin{frame}{Cryptography and Hardness Assumptions}
\begin{itemize}
    \item Modern cryptography exploits asymmetry between easy verification and hard inversion.
    \item Examples:
    \begin{itemize}
        \item integer factorization,
        \item discrete logarithms,
        \item elliptic-curve discrete logarithms,
        \item cryptographic hash functions.
    \end{itemize}
    \item These are evidence that hardness matters in practice.
    \item But they do \textbf{not} directly resolve $\mathsf{P}$ vs. $\mathsf{NP}$.
\end{itemize}
\end{frame}

\begin{frame}{Blockchain and Post-Quantum Context}
\begin{itemize}
    \item Blockchain systems test cryptographic hardness at global scale under adversarial conditions.
    \item Zero-knowledge systems highlight the distinction between efficient verification and expensive discovery.
    \item Quantum computing increases pressure on classical hardness assumptions.
    \item Post-quantum cryptography shifts attention to lattices, hash-based signatures, and new security assumptions.
\end{itemize}
\end{frame}

\section{AI Co-Pilot Framework}

\begin{frame}{Why an AI Co-Pilot?}
The claim is \textbf{not} that AI proves or disproves the conjecture.

\medskip
The proposal is methodological:
\begin{itemize}
    \item organize known results,
    \item map dependencies among reductions and lemmas,
    \item generate constrained conjectures,
    \item explore structured instances empirically,
    \item interface with theorem provers and symbolic systems.
\end{itemize}

\medskip
Workflow:
\begin{center}
\emph{human judgment + AI-assisted search + formal/computational validation}
\end{center}
\end{frame}

\begin{frame}{Research Directions for the Framework}
\begin{itemize}
    \item \textbf{Restricted classes:} bounded treewidth, sparse SAT, planted distributions.
    \item \textbf{Reduction mining:} catalog shared structural bottlenecks among NP-complete problems.
    \item \textbf{Proof-barrier cartography:} classify which arguments relativize, algebrize, or resemble natural proofs.
    \item \textbf{Cryptographic interface:} connect hardness assumptions, pseudorandomness, and lower-bound barriers.
\end{itemize}
\end{frame}

\begin{frame}{Known Meta-Barriers}
Any successful proof likely has to escape several known barriers:
\begin{itemize}
    \item \textbf{Relativization}
    \item \textbf{Natural Proofs}
    \item \textbf{Algebrization}
\end{itemize}

\medskip
These barriers help explain why the problem has resisted decades of work.
\end{frame}

\section{Takeaways}

\begin{frame}{Key Takeaways}
\begin{itemize}
    \item P vs\@. NP asks whether efficient verification implies efficient solution.
    \item The core practical issue is resource-bounded feasibility.
    \item NP-complete problems model combinatorial explosion.
    \item Cryptography and large-scale systems make hardness assumptions concrete.
    \item AI may improve the \emph{research process}, even if the proof itself still needs new mathematics.
\end{itemize}
\end{frame}

\begin{frame}{Conclusion}
\begin{itemize}
    \item The P vs\@. NP problem remains one of the deepest open questions in mathematics and theoretical computer science.
    \item Its significance extends far beyond abstract definitions: it shapes optimization, cryptography, verification, and the theory of computation.
    \item An AI co-pilot framework is best viewed as a disciplined assistant for exploration, synthesis, and formal checking.
    \item The conjecture itself still demands genuinely new mathematical insight.
\end{itemize}
\end{frame}

\begin{frame}[plain]
\centering
\Huge Questions?
\end{frame}

\end{document}
